% Shared options for presentation slides.
\ExplSyntaxOn
\keys_define:nn { presentationslide }
    {
        title .tl_set:N = \slidetitle,
        subtitle .tl_set:N = \slidesubtitle,
        note .tl_set:N = \slidenote,
        image .tl_set:N = \slideimage,
        image1 .tl_set:N = \slideimageone,
        image2 .tl_set:N = \slideimagetwo,
        caption .tl_set:N = \slidecaption,
        caption1 .tl_set:N = \slidecaptionone,
        caption2 .tl_set:N = \slidecaptiontwo,
        source .tl_set:N = \slidesource,
        source1 .tl_set:N = \slidesourceone,
        source2 .tl_set:N = \slidesourcetwo,
        height .tl_set:N = \slideheight
    }
\NewDocumentCommand{\slidesetup}{m}
    { \keys_set:nn { presentationslide } { #1 } }
\ExplSyntaxOff

% Display an optional source below an image or table.
\newcommand{\showslidesource}[1]{%
    \ifx#1\empty\else
        {\footnotesize\sourcelabel: #1\par}%
    \fi
}

% Use the current section, subsection, or subsubsection when no title is supplied.
\newcommand{\slideheading}{%
    \ifx\slidetitle\empty
        \ifx\insertsubsubsectionhead\empty
            \ifx\insertsubsectionhead\empty
                \frametitle{\insertsectionhead}
            \else
                \global\presentationtitlesubsectiontrue
                \frametitle{\insertsubsectionhead}
            \fi
        \else
            \frametitle{\insertsubsubsectionhead}
        \fi
    \else
        \frametitle{\slidetitle}
    \fi
    \framesubtitle{\slidesubtitle}%
}

% Content slides share a heading and an optional speaker note.
\newcommand{\slideframe}[2]{%
    \slidesetup{#1}%
    \begin{frame}
        \slideheading
        #2
        \note{\slidenote}
    \end{frame}%
}

% \presentationimage{file}{height}{maximum width}
\newcommand{\presentationimage}[3]{%
    \includegraphics[width=#3,height=#2,keepaspectratio]{#1}%
}

% Bullet list beside an image.
\NewDocumentEnvironment{BulletImageSlide}{O{} +b}{%
    \slideframe{height=4.2cm, #1}{%
        \begin{columns}[c,onlytextwidth]
            \begin{column}{0.50\textwidth}
                \begin{itemize}#2\end{itemize}
            \end{column}
            \begin{column}{0.46\textwidth}
                \begin{figure}
                    \centering
                    \includegraphics[width=\linewidth,height=\slideheight,keepaspectratio]{\slideimage}
                    \caption{\slidecaption}
                    \showslidesource{\slidesource}
                \end{figure}
            \end{column}
        \end{columns}
    }%
}{}

% Title slide with an optional speaker note.
\NewDocumentCommand{\TitleSlide}{O{}}{%
    \begingroup
    \slidesetup{#1}%
    \begin{frame}[plain,noframenumbering]
        \titlepage
        \note{\slidenote}
    \end{frame}
    \endgroup
}

% Table of contents with an optional subtitle and speaker note.
\NewDocumentCommand{\ContentsSlide}{O{}}{%
    \begingroup
    \slidesetup{#1}%
    \begin{frame}{\contentsname}{\slidesubtitle}
        \tableofcontents[subsectionstyle=show]
        \note{\slidenote}
    \end{frame}
    \endgroup
}

% A single image scaled to its height and cropped to the available width.
\NewDocumentCommand{\ImageSlide}{O{}}{%
    \begingroup
    \slideframe{height=5.2cm, #1}{%
        \begin{figure}
            \centering
            \presentationimage{\slideimage}{\slideheight}{0.95\linewidth}
            \caption{\slidecaption}
            \showslidesource{\slidesource}
        \end{figure}
    }%
    \endgroup
}

% One panel of an image pair.
\newcommand{\slideimagepanel}[3]{%
    \begin{subfigure}[t]{0.46\linewidth}
        \vspace{0pt}\centering
        \presentationimage{#1}{\slideheight}{\linewidth}
        \caption{#2}
        \showslidesource{#3}
    \end{subfigure}%
}

% Two images with individual captions and an optional shared caption.
\NewDocumentCommand{\ImagePairSlide}{O{}}{%
    \begingroup
    \slideframe{height=3.8cm, #1}{%
        \begin{figure}\centering
            \slideimagepanel{\slideimageone}{\slidecaptionone}{\slidesourceone}%
            \hfill
            \slideimagepanel{\slideimagetwo}{\slidecaptiontwo}{\slidesourcetwo}%
            \ifx\slidecaption\empty\else\caption{\slidecaption}\fi
            \showslidesource{\slidesource}
        \end{figure}
    }%
    \endgroup
}

% Centered content with an optional title, subtitle, and speaker note.
\NewDocumentEnvironment{EmphasisSlide}{O{} +b}{%
    \slideframe{#1}{%
        \centering
        #2\par
    }%
}{}

% References with automatic continuation across slides.
\NewDocumentCommand{\ReferencesSlide}{O{}}{%
    \begingroup
    \slidesetup{#1}%
    \setbeamertemplate{frametitle continuation}[from second]
    \begin{frame}[allowframebreaks]{\refname}{\slidesubtitle}
        \printbibliography[heading=none]
        \note{\slidenote}
    \end{frame}
    \endgroup
}

% A table with a caption and freely formatted rows.
\NewDocumentEnvironment{TableSlide}{O{} +b}{%
    \slideframe{#1}{%
        \begin{table}\centering
            \caption{\slidecaption}
            {\usebeamerfont{table body}#2\par}
            \showslidesource{\slidesource}
        \end{table}
    }%
}{}

% A slide with a list of text.
\NewDocumentEnvironment{TextSlide}{O{} +b}{%
    \slideframe{#1}{%
        \begin{itemize}
            \setlength\itemsep{3ex}
            #2
        \end{itemize}
    }%
}{}

% A fragile frame for verbatim code entered directly in the slide.
\newenvironment{CodeSlide}[1][]{%
    \slidesetup{#1}%
    \begin{frame}[fragile,environment=CodeSlide]
        \slideheading
        \begin{center}\begin{minipage}{0.46\textwidth}
}{%
        \end{minipage}\end{center}
        \note{\slidenote}
    \end{frame}
}

% A slide with freely formatted formulas and legends.
\NewDocumentEnvironment{FormulaSlide}{O{} +b}{%
    \slideframe{#1}{{\usebeamerfont{slide body}#2\par}}%
}{}
